\chapter{Expanding the Experiential Possibilities of Biophilic Interaction}
\label{ch:biophilic-futures}

\begin{synopsis}
\synopsislead{What forms of indoor experience become possible}{when nature is approached as more than a visual or restorative resource? The preceding chapter established an affective approach for investigating how indoor conditions acquire significance through bodily, spatial, temporal, social, and interpretive relations. This chapter turns from investigating existing experience towards considering what richer forms of experience future interactive indoor environments might support.}

This chapter develops this direction through multisensory biophilic design. Drawing on semi-structured interviews with 13 experts across architecture, design, and research, it examines what the legacy, practice, and imaginaries of biophilic architecture can contribute to interaction design. The analysis develops seven topics describing the breadth of expert discourse, eight themes articulating expert imaginaries of biophilic futures, and five design opportunities for technologically mediated indoor environments.

Within the trajectory of the dissertation, this chapter develops the first part of \rqref{biophilic} by expanding the experiential scope of biophilic interaction. The findings position nature as ecological process, sensory condition, cultural relation, and more-than-human presence rather than principally as visual greenery. They draw attention to multisensory and bodily engagement, natural variability, ecological authenticity, (dis)comfort, equity, and relationships of care and responsibility towards the living world.

These findings establish biophilic design as a productive approach for imagining experiential possibilities beyond prevailing models of indoor comfort and optimisation. They also create a further design problem for the dissertation: how can these possibilities be translated from a broad set of architectural imaginaries into concepts that can more systematically support interaction design?
\end{synopsis}

\chaptersource{%
This chapter is based on: Shruti Rao, Judith Good, and Hamed Alavi.
\emph{Designing Multisensory Biophilic Futures: Exploring the Potential of Interaction Design to Deepen Human Connections with Nature in Indoor Environments}.
Presented at DIS 2025, Funchal, Portugal.%
}

\clearpage

\import{papers/04-bidf/}{thesis-body.tex}
\clearpage

\begin{chapterclosing}{Concluding synthesis}
This chapter broadens the experiential and ecological possibilities considered within biophilic interaction. The expert accounts show that meaningful relations with nature indoors can involve sensory variation, bodily movement, ecological processes, cultural interpretation, multispecies presence, discomfort, and forms of care and responsibility that extend beyond the provision of greenery or restorative visual experience.

The five design opportunities translate these perspectives into directions for interactive indoor environments: supporting emergent natural processes, attending to bodily sensitivity, engaging carefully with biophilic (dis)comfort, considering equitable access to nature-rich experience, and supporting relational encounters with natural processes. Together, they position multisensory and multidimensional experience as a basis for considering what technologically mediated biophilic environments might enable. The paper itself frames these contributions as an expansion of HBI's experiential and ecological vocabulary. 

Within the dissertation, this chapter therefore establishes the experiential scope of the biophilic design direction. It identifies a set of possibilities and values drawn from architectural knowledge, while leaving open how these possibilities can be organised into a more explicit interaction design resource.

\chaptercarryforward{The following chapter moves from expanding the possibilities of biophilic interaction to structuring them for design, using expert design sessions to articulate what can be configured, why particular experiential qualities matter, and how interaction with biophilic indoor environments might occur.}
\end{chapterclosing}